\chapter{概率空间与期望}\label{chap:prob}

本章用最少的记号建立概率论的语言。English and 中文 are mixed on purpose so that
hover, completion and live preview see realistic text around the formulas.

\section{基本定义}\label{sec:prob-basics}

\begin{definition}{概率空间 Probability space}{prob-space}
  一个\term{概率空间}是三元组 $(\Omega, \sigmaF, \Prob)$，其中 $\Omega$ 是样本空间，
  $\sigmaF$ 是 $\Omega$ 上的 $\sigma$-代数，$\Prob\colon \sigmaF \to [0,1]$ 满足
  $\Prob(\Omega) = 1$ 且对两两不交的 $A_1, A_2, \dots \in \sigmaF$ 有
  \[
    \Prob\Bigl(\bigcup_{i=1}^{\infty} A_i\Bigr) = \sum_{i=1}^{\infty} \Prob(A_i).
  \]
\end{definition}

随机变量 $X\colon \Omega \to \R$ 的期望写作 $\E{X} = \int_\Omega X \dd\Prob$；
换一个测度 $Q$ 时写作 $\E[Q]{X}$。若 $X_1, \dots, X_n \iid P$，样本均值为
\(\bar X_n = \frac{1}{n}\sum_{i=1}^n X_i\)。

示性函数的定义用 \texttt{cases} 环境：
\begin{equation}\label{eq:indicator}
  \mathbf{1}_A(\omega) =
  \begin{cases}
    1, & \omega \in A, \\
    0, & \text{否则 (otherwise)}.
  \end{cases}
\end{equation}

\begin{note}
  式~\eqref{eq:indicator} 的期望就是概率：$\E{\mathbf{1}_A} = \Prob(A)$。
  价格写成 \$5 时美元符号被转义，不是数学；\verb|$x$| 里的美元符号也不是数学。
\end{note}

% 注释里的公式不应被渲染：$\alpha + \beta$ \[ \int f \]

\section{方差与条件期望}\label{sec:variance}

方差有两种等价写法：
\begin{align}
  \Var(X) &= \E{(X - \E{X})^2} \label{eq:var-def} \\
          &= \E{X^2} - \bigl(\E{X}\bigr)^2. \label{eq:var-short}
\end{align}
协方差 $\Cov(X, Y) = \E{XY} - \E{X}\E{Y}$，且 $\Var(X + Y) = \Var(X) + \Var(Y) + 2\Cov(X, Y)$。

\begin{theorem}{全期望公式 Law of total expectation}{total-exp}
  若 $\E{\abs{X}} < \infty$，则
  \begin{equation}\label{eq:total-exp}
    \E{X} = \E{\E{X \given Y}}.
  \end{equation}
\end{theorem}

\begin{proof}
  对离散的 $Y$，
  \begin{align*}
    \E{\E{X \given Y}}
      &= \sum_{y} \Prob(Y = y)\, \E{X \given Y = y} \\
      &= \sum_{y} \sum_{x} x\, \Prob(X = x, Y = y)
       = \E{X}.
  \end{align*}
  一般情形由单调类定理得到，细节见~\cite[第 2 章]{zhang2020notes}。
\end{proof}

\begin{example}[抛硬币 Coin tossing]
  设 $X \sim \mathrm{Bernoulli}(p)$，则 $\E{X} = p$，由~\eqref{eq:var-short} 得
  $\Var(X) = p - p^2 = p(1-p)$。取 $p = \tfrac12$ 时方差最大。
\end{example}

\begin{remark}
  定理~\ref{thm:total-exp} 在第~\ref{chap:linalg} 章会以投影的形式重新出现：
  条件期望是 $L^2$ 中到子空间的正交投影，参见命题~\ref{pro:psd} 和图~\ref{fig:projection}。
\end{remark}

\section{集中不等式}\label{sec:concentration}

常用的几个不等式：
\begin{enumerate}[(a)]
  \item Markov：对 $a > 0$，$\Prob(\abs{X} \ge a) \le \E{\abs{X}} / a$；
  \item Chebyshev：$\Prob(\abs{X - \E{X}} \ge \eps) \le \Var(X)/\eps^2$；
  \item Hoeffding：若 $X_i \in [a_i, b_i]$ 独立，则
    \begin{itemize}
      \item 单侧界 $\Prob(\bar X_n - \E{\bar X_n} \ge t) \le \exp\!\left(-\frac{2n^2t^2}{\sum_i (b_i - a_i)^2}\right)$；
      \item 双侧界再乘以 $2$。
    \end{itemize}
\end{enumerate}

\begin{description}
  \item[弱大数定律] $\bar X_n \to \E{X_1}$ 依概率收敛；
  \item[强大数定律] 几乎必然收敛，见~\cite{li2019lln,smith2021concentration}。
\end{description}

Plain \TeX{} display math is deprecated but common in real notes:
$$
  \KL{P}{Q} = \int \log\frac{\dd P}{\dd Q} \dd P \ge 0,
$$
and the forward divergence with an optional argument is $\KL[\text{fwd}]{P}{Q}$.%
\footnote{脚注里的数学：$\KL{P}{P} = 0$。}
